% Chapter 9, Section 4

\section{Applications of CNNs}
\label{sec:cnn-applications}

\subsection*{Intuition}
Backbones of stacked convolutions extract spatially local features that become increasingly abstract with depth. Task-specific heads (classification, detection, segmentation) transform backbone features into outputs appropriate to the problem \cite{GoodfellowEtAl2016,Prince2023}.
\subsection{Image Classification}

Task: Assign label to entire image.

Architecture:
\begin{itemize}
    \item Convolutional layers extract features
    \item Pooling reduces dimensions
    \item Fully connected layers for classification
    \item Softmax output
\end{itemize}

\subsection{Object Detection}

Task: Localize and classify objects in images.

\textbf{Region-based methods (R-CNN family):}
\begin{itemize}
    \item R-CNN: region proposals + CNN features
    \item Fast R-CNN: end-to-end training
    \item Faster R-CNN: learned region proposals
    \item Mask R-CNN: adds instance segmentation
\end{itemize}

\textbf{Single-shot methods:}
\begin{itemize}
    \item YOLO (You Only Look Once): real-time detection
    \item SSD (Single Shot Detector): multi-scale detection
\end{itemize}

\subsection{Semantic Segmentation}

Task: Assign class label to each pixel. U-Net popularized encoder-decoder with skip connections for biomedical images \cite{Ronneberger2015}.

\textbf{FCN (Fully Convolutional Networks):}
\begin{itemize}
    \item Replace FC layers with convolutions
\end{itemize}
